\section{Conclusiones y Trabajos Futuros}

\subsection{Conclusiones principales}

Particularmente llamativo es el hecho de que una plataforma consolidada como el VRSS, diseñada hace más de una década bajo paradigmas de observación terrestre tradicionales \cite{VRSS1_2012}, pueda extender significativamente su utilidad mediante una capa complementaria de nanosatélites. Los resultados obtenidos demuestran de forma consistente que la arquitectura híbrida propuesta reduce la revisita efectiva de varios días a pocas horas, eleva la continuidad temporal por encima del 90\% y mantiene la alta resolución espacial inherente al segmento principal.

Las simulaciones confirman mejoras cuantificables en latencia y capacidad de respuesta ante eventos dinámicos. Más allá de los números, esta integración transforma el sistema de una herramienta de cartografía periódica en una infraestructura de inteligencia ambiental con verdadero valor operativo. La aproximación no busca reemplazar lo existente, sino potenciarlo de manera inteligente, aprovechando las fortalezas de cada componente.

\subsection{Recomendaciones}

Uno de los desafíos persistentes que surge al cerrar este tipo de estudios radica en traducir hallazgos técnicos en acciones concretas. Se recomienda priorizar el desarrollo de un prototipo de tres a cinco nanosatélites para validar la integración en órbita real y refinar los protocolos de fusión y tasking coordinado. 

Resulta igualmente importante establecer marcos de colaboración interinstitucional que involucren a la Agencia Bolivariana para Actividades Espaciales, universidades y posibles socios internacionales. Desde el punto de vista técnico, invertir en capacidades locales de calibración cruzada y procesamiento en borde se presenta como paso crítico para reducir dependencia de infraestructura externa.

En el plano operativo, las agencias usuarias deben comenzar desde ya la actualización de sus flujos de trabajo para incorporar productos de alta frecuencia temporal. Ignorar este aspecto podría limitar seriamente el impacto real del sistema híbrido.

\subsection{Trabajos futuros}

Lejos de ser un asunto resuelto, el camino hacia sistemas de observación terrestre verdaderamente en tiempo real sigue abierto. Una línea prometedora consiste en incorporar enlaces inter-satélite (crosslinks) y sistemas de propulsión de baja potencia que permitan reconfiguración orbital dinámica y reducción adicional de latencia \cite{Chan2024}.

Entre las extensiones naturales se encuentran: el desarrollo de algoritmos de inteligencia artificial embarcados para detección autónoma de eventos, la integración con datos de otras constelaciones internacionales bajo esquemas de acceso compartido, y la exploración de sensores adicionales (SAR compacto o hiperespectral) en la capa de nanosatélites. 

También merece atención el análisis de viabilidad a mayor escala, incluyendo modelos de negocio para sostenibilidad financiera y evaluación de impacto socioeconómico en aplicaciones prioritarias para Venezuela. 

En definitiva, la presente investigación establece una base sólida sobre la cual construir capacidades espaciales más resilientes y oportunas. Los próximos pasos determinarán si esta visión híbrida se convierte en un referente regional para países en desarrollo con aspiraciones de soberanía tecnológica.